\documentclass[10pt,a4paper]{amsart}
\usepackage[margin=19mm]{geometry}
\usepackage{amsmath,amssymb,mathtools}
\usepackage[scr=rsfs,cal=euler]{mathalfa}
\usepackage{xfrac}
\usepackage[unicode,hidelinks,bookmarks=false]{hyperref}
\newcommand{\ie}{i.e.\ }
\newcommand{\cf}{cf.\ }
\newcommand{\resp}{respectively\ }
\newcommand{\Subsection}{\subsection}
\providecommand{\nobreakdash}{\nobreak\hbox{-}}
\setlength{\emergencystretch}{2em}
\multlinegap0pt
\numberwithin{equation}{section}
\newtheorem{theo}{Theorem}[section]
\newtheorem{theorem}[theo]{Theorem}
\newtheorem{lemma}[theo]{Lemma}
\newtheorem{definition}[theo]{Definition}
\newtheorem{proposition}[theo]{Proposition}
\newtheorem{remark}[theo]{Remark}
\newtheorem{corollary}[theo]{Corollary}
\newtheorem{example}[theo]{Example}
\allowdisplaybreaks
\title[Pricing callable bonds and optimal callable time under the Fractional Black-Scholes market]{Pricing callable bonds and optimal callable time under the fractional Black-Scholes market: L1 consistency of the deep-learning price estimator (Theorem 3.6)}
\author{Yuecai Han, Yinong Wu, Xudong Zheng}
\date{Source: Electronic Journal of Differential Equations 2026, No. 09, pp. 1--13; DOI 10.58997/ejde.2026.09; publisher version}
\begin{document}
\maketitle
\noindent\emph{Electronic Journal of Differential Equations}, Vol. 2026 (2026), No. 09, pp. 1--13.\newline
ISSN: 1072-6691. DOI \href{https://doi.org/10.58997/ejde.2026.09}{10.58997/ejde.2026.09}. This work is licensed under a CC BY 4.0 license.\par\smallskip
\noindent\textit{Editorial scope.} Counted claim: original Theorem 3.6 (the main consistency theorem of the article; source lines 853-858) together with its complete original proof at source lines 860-896, which already contains the proof environment, delivered as the single primary proof body. The whole of Sections 1 to 3 is retained uncounted as complete context (source lines 60-851 and 898-906): the callable-bond pricing problem, the fractional Brownian motion and Wick product setting, the fractional Black-Scholes discretisation, Assumptions (A1) and (A2), the Glivenko-Cantelli type uniform convergence of the conditional distribution estimator (Lemma 3.1), the boundedness-based L1 convergence of the continuation-value estimator (Lemma 3.2), the backward-recursion stability estimate (Lemma 3.3), the truncation argument that removes the boundedness hypothesis (Lemma 3.4), the convergence of the estimated optimal call time (Lemma 3.5) and Remark 3.7 on the dual upper bound. Section 4 (numerical simulations) and Section 5 (conclusions), the acknowledgements and the bibliography are not selected and are not used to inflate H. Because no content line of Sections 1 to 3 is dropped, every printed equation, lemma, theorem and remark number of the delivered text coincides with the published PDF, in particular Lemma 3.1 to Lemma 3.5, Theorem 3.6 and Remark 3.7. One printed cross-reference whose target lies outside the delivered unit, namely the roadmap sentence that announces the numerical-simulation and conclusion sections, is delivered as a hyperlink whose visible text is the published section number. Slips kept literal and recorded without mathematical repair: in the counted proof the upper summation index of the empirical average is printed as k\_N-1 instead of the braced subscript k\_\{N-1\} used everywhere else, and the estimator is printed as V-hat without its argument while the surrounding display carries the argument r\_0; both printed forms are reproduced verbatim. No mathematics is written, repaired or re-derived; all mathematics is native text with no image dependency.
\section{Introduction}

The valuation of callable bonds presents a significant challenge in financial 
engineering due to their embedded option features. The models estimate the bond's 
value based on possible callable paths, assigning probabilities to different scenarios. 
This approach renders the bond price highly sensitive to the timing of the call, 
as early or delayed call can significantly affect the bond's ultimate payoff and, 
therefore, its market value.
The optimal redemption time for a callable bond is influenced by interest rates, 
time to maturity, and potential call features. This timing directly impacts callable 
bond pricing models, which aim to incorporate optimal callable strategies to determine 
fair value.

For investigating callable bonds, Narayanan and Lim \cite{narayanan1989call} explored 
the optimality of calling corporate zero-coupon bonds for refunding, considering 
corporate tax effects. Their analysis reveals that refunding is not optimal unless the 
corporate tax rate exceeds $50\%$. The study also shows that callable zero-coupon bonds 
more frequently include restrictive covenants, indicating the call feature offers firms 
flexible, low-cost options for future recapitalization.
Farto and V\'azquez \cite{farto2005numerical} focused on a complex financial 
product==a coupon-bearing callable bond with a compulsory notice period==by 
addressing the price discontinuities caused by discrete coupons and call features. 
The bond's value depends on time and a stochastic interest rate, modeled through a
Black-Scholes type PDE with appropriate boundary and jump conditions. 
To solve this, the authors proposed a numerical method combining characteristics-based 
time discretization with finite element methods for the interest rate variable. 
The method shows strong performance when tested against known solutions under the 
Vasicek and Cox-Ingersoll-Ross (CIR) models, and outperforms prior finite volume 
approaches in the call-with-notice scenario.

Aron \cite{aaron2021zero} estimated a dynamic term structure model directly on treasury 
coupon bond data, bypassing the need for a zero-coupon yield curve. 
A linearity generating approach enables fast estimation and separates cross-sectional 
from time series parameters. The model quantifies the on-the-run and “notes versus 
bonds” premiums from 1990 to 2017 within a unified, no-arbitrage framework.
D{\'\i}az et al.\ \cite{diaz2022yield} compared various zero-coupon yield curve datasets,
which differ by asset selection and fitting methods, despite representing the same 
underlying reality. Using an empirical analysis on callable bond pricing, 
it finds that results can vary significantly depending on the dataset used, 
particularly due to differences in volatility inputs. This sensitivity may create 
incentives for financial institutions to selectively choose datasets in their favor, 
posing a potential moral hazard.

Lin and Zhu \cite{lin2022pricing} addressed the pricing of callable-puttable 
convertible bonds using an integral equation (IE) approach. The complexity arises 
from the interplay of call, put, and conversion features, leading to scenarios with 
one or two moving boundaries depending on time to maturity and contract terms. 
The analysis identifies critical time points where the structure of the problem changes, 
requiring different pricing strategies==ranging from handling only the put feature 
to managing both put and conversion, or reducing to a simpler vanilla bond case. 
Each scenario necessitates a distinct numerical solution of corresponding IE systems.

Skalicky et al.\ \cite{skalicky2022valuation} presented a new method for valuing the 
risk of bond buybacks by issuers at a specified call price, specifically for privately 
traded bonds with embedded European or multiple options. Unlike traditional models 
for marketable callable bonds, this approach addresses the challenges of limited 
transaction data. Its modular structure enables the incorporation of issuer behavior, 
alternative investment opportunities, and varying interest rate expectations.

 Kusnadi et al.\ \cite{kusnadi2023callable} focused on the impact of early redemption
and default risk on the value of Indonesian bonds using binomial interest rate tree 
models. Simulations reveal that early redemption risk generally lowers a bond's present
value, while default risk, under certain assumptions, can unexpectedly increase it. 
The higher bond value under default risk is attributed to an inflated recovery 
fraction in the first period, driven by data limitations.
Dow and Orfanos \cite{dow2025interest} addressed the limitations of duration and 
convexity matching in hedging fixed income securities with negative convexity, 
such as callable bonds and mortgage-backed securities. It introduces the concepts of 
bond tilt and bond agility to manage residual risk, deriving approximation formulas 
that incorporate higher-order effects. The proposed methods closely track price-yield 
dynamics, achieving mean absolute errors below 2.5$\%$ across various callable bonds 
under yield shifts of up to $\pm 200$ basis points.


 Hobson et al.\ \cite{hobson2024callable} examined the callable convertible bond 
problem under a liquidity constraint represented by Poisson signals. Unlike traditional 
models, it assumes that when the bondholder and the firm act simultaneously, neither 
has full priority, instead, a proportion $ m \in [0, 1] $ of the bond is converted to 
equity, while the remainder is called by the firm. This framework extends the special 
case analyzed by Liang and Sun \cite{liang2019dynkin}, where the bondholder held full 
priority ($ m = 1 $), and provides a comprehensive solution to the problem under 
this more general setting.

In addition, some scholars use deep learning methods to deal with option and bond 
pricing.
Becker et al.\ \cite{becker2019deep} proposed a deep learning method for optimal 
stopping problems which directly learns the optimal stopping rule from Monte Carlo 
samples.
 Accurate results are obtained for the problem of optimal stopping
a fractional Brownian motion with short computing times.
Becker et al.\ \cite{becker2021solving} presented a deep learning-based algorithm 
to price early exercise options, effectively addressing high-dimensional optimal 
stopping problems common in derivatives such as American and Bermudan options. 
The method approximates both the option price and optimal exercise strategy and can be 
applied to other optimal stopping problems with simulatable stochastic processes. 
Numerical results, including cases with up to 5000 dimensions, demonstrate the 
algorithm's accuracy and efficiency compared to existing benchmarks.

Tan et al.\ \cite{tan2022deeppricing} introduced Deeppricing model, a data-driven model 
for pricing convertible bonds using a novel financial time series GAN called FinGAN. 
FinGAN accurately captures complex features of stock return dynamics, such as fat tails, 
long-range dependence, and asymmetry, and transitions to a risk neutral distribution, 
enabling more realistic and flexible pricing. Experiments on the Chinese convertible 
bond market show that Deeppricing model outperforms traditional models (e.g., 
Black-Scholes, CEV, GARCH) and other GAN-based methods, especially for higher 
volatility or equity-like bonds. %In addition, investment strategies based on 
Deeppricing model yield strong annualized returns, demonstrating their practical value.

The framework of this paper is as follows. 
Section \ref{sec:basics} introduces the pricing problem of callable bonds. 
We apply deep learning method to calculate the lower bound and construct confidence 
intervals for the price of callable bonds in Section \ref{sec:main}. 
The results of the numerical simulation are presented in Section \href{https://ejde.math.txstate.edu/Volumes/2026/09/han.pdf}{4}.
Conclusions and future work are drawn in Section \href{https://ejde.math.txstate.edu/Volumes/2026/09/han.pdf}{5}.


\section{Problem formulation and fractional Black-Scholes model}
\label{sec:basics}

\subsection{Problem formulation}
\label{lower}
Zero-coupon callable bond pricing problem can be formulated through an interest rate 
process
$ r_{t}$, $0 \leq t \leq T $, defined on $\left( \Omega, \mathcal{F}, P \right)$, where 
$P$ is the risk-neutral probability measure. The set of stopping times considered in 
this paper is finite. During the life of the bond, the issuer has a finite number of 
time points $t_0=0,t_1,\dots,t_N=T$, to choose whether to call the bond or not.
If called at time $t_{i}$, $i=0,1,2, \dots, N$, the issuer has the right to call the
 bond by paying a predetermined price $S$ to bond holder.
The payoff functions
$p_{t_i}\left(r_{t_{i}}\right.$) are some known functions 
$p_{t_i}(\cdot):R \to [0,+\infty)$.
In this paper, we set the payoff function is
$p_{t_i}(t_i)=\min \big(K, F \cdot e^{-\int_0^{t_i} r(s) d s}\big)$.
If the payoff at $t_i$ is lower than the continuation value at that time, the 
issuer early call the bond; otherwise, continue to issue it. The bond values and the
continuation value satisfy
\begin{gather*}
V_{t_i}=\sup _{\tau \in \mathcal{T}_{i}} \mathbb{E}\left[D\left(t_i, \tau\right) 
p_{\tau}\left(r_{\tau}\right) \mid \mathcal{F}_{t_i} \right]
=\mathbb{E}\left[D\left(t_i, \tau^{\ast }_i \right) p_{\tau^{\ast }_i }
\left(r_{\tau^{\ast }_i }\right) \mid \mathcal{F}_{t_i} \right],
\\
 C_{t_i}=\mathbb{E}\left[D\left(t_i, t_i+1\right) V_{t_i+1}
 \left(r_{t_i+1}\right) \mid \mathcal{F}_{t_i}\right],
\end{gather*}
where $\mathcal{T}_{i}=\left\{t_{i}, t_{i+1}, \dots, t_{N}\right\}$,
$\tau^{\ast }_i$ is the optimal moment to call, 
$D\left(t, \tau\right)=e^{-\int_{t}^{\tau} r_{u} d u}$ is the discount factor,
$r_{u}$ is the interest rate process we analyze through the paper.
For simplicity, we let $\mathcal{F}_i$ denote $\mathcal{F}_{t_i}$, 
$p_{i}\left(r_{i}\right)$ denote $p_{t_{i}}\left(r_{t_{i}}\right)$, $V_{i}$ denote $V_{t_i}$, and $C_{i}$ denote $C_{t_i}$. Then, the bond values and the continuation values satisfy the dynamic programming equations
\begin{equation}
\label{dynamic}
\begin{gathered}
V_{N}=p_{N}\left(r_N\right), \\
V_{i}=\max \Big\{p_{i}\left(r_i\right), \mathbb{E}\Big[
e^{-\int_{t_i}^{t_{i+1}} r_{u} d u} V_{i+1}\left(r_{i+1}\right) 
\mid \mathcal{F}_i\Big]\Big\}, \quad i=0,1, \dots, N-1 ,\\
C_{N}=0, \\
C_{i}=\mathbb{E}\Big[
e^{-\int_{t_i}^{t_{i+1}} r_{u} d u}
\max \left\{p_{i+1}\left(r_{i+1}\right), C_{i+1}
\left(r_{i+1}\right)\right\} \mid \mathcal{F}_i\Big], \quad i=0,1, \dots, N-1 ,
\end{gathered}
\end{equation}
and the bond values satisfy
\[
V_{i}=\max \left\{p_{i}\left(r_i\right), C_{i}\left(r_i\right)\right\} .
\]

\subsection{Fractional Brownian motion and Wick product}
The fractional Brownian motion with Hurst parameter $H\in\left( 0,1\right) $ is 
a zero-mean Gaussian process $\left ( B_{t}^{H} \right )_{t \in [0,\infty)} $ with
 covariance
\[
\mathbb{E}\left[B_{t}^{H} B_{s}^{H}\right]=\frac{1}{2}\left(t^{2 H}+s^{2 H}
-|t-s|^{2 H}\right), \quad s,t \in [0,\infty).
\]
The fractional Brownian motion is self-similar, i.e. $ B_{at} $ and 
$a^{H} B_{t}^{H}$ have the same probability law for all $a>0$. Compared to 
Brownian motion, it displays positive correlation property when $\frac{1}{2}<H<1$
and negative correlation property when $0<H<\frac{1}{2}$. This special property 
of fractional Brownian motion allows it to describe models with long-range dependence. The fractional Brownian increment
\[
 \Delta B_{t, s}^H=B_t^H-B_s^H, \quad 0\leq s<t \leq\infty,
\]
has the following properties
\begin{gather*}
\mathbb{E}\left[\Delta B_{t, s}^H\right]
=\mathbb{E}\left[B_t^H-B_s^H\right]=0,\\
\mathbb{E}\left[\left(\Delta B_{t, s}^H\right)^2\right]
=\mathbb{E}\left[\left(B_t^H-B_s^H\right)\left(B_t^H-B_s^H\right)\right]=|t-s|^{2 H} ,\\
\mathbb{E}\left[\Delta B_{t, s}^H \Delta B_{s, 0}^H\right]
=\mathbb{E}\left[\left(B_t^H-B_s^H\right)\left(B_s^H-B_0^H\right)\right] =\frac{1}{2}\left[t^{2 H}-s^{2 H}-(t-s)^{2 H}\right] .
\end{gather*}

In this paper, we consider the case $\frac{1}{2}<H<1$. We use the following 
stochastic integral representation of the fractional Brownian motion
\[
B_t^H=\int_0^t z_H(t, s) d B_s, \quad t \in [0,\infty),
\]
where $B_{t} $ is a standard Brownian motion, and the deterministic kernel
\[
 z_H(t, s)=1_{\{t \geq s\}} c_H\big(H-\frac{1}{2}\big) s^{\frac{1}{2}-H} 
 \int_s^t u^{H-\frac{1}{2}}(u-s)^{H-\frac{3}{2}} d u
\]
with the constant
\[
c_H=\sqrt{\frac{2 H \Gamma\left(\frac{3}{2}-H\right)}{\Gamma\left(H+\frac{1}{2}\right) 
\Gamma(2-2 H)}} ,
\]
where $\Gamma$ is the Gamma function. Because of the existence of a bijective transfer 
operator, the filtration generated by $B_{t}^{H}$ is also the one generated by 
$B_{t}$ through expressing $B_{t}$ as a Wiener integral with respect to $B_{t}^{H}$.

The integral theory of fractional Brownian motion based on the Wick product is 
introduced by \cite{duncan2000stochastic} and further developed by 
\cite{hu2003fractional} through fractional white noise theory, which proves that 
the financial markets are arbitrage free and complete when the integrals defined by 
this method are applied to financial market models. Define the fractional kernel 
$\phi: \mathbb{R}^2 \to \mathbb{R}$, the norm $\|\cdot\|_\phi^2$ and inner
product $\langle\cdot, \cdot\rangle_\phi$ by
\begin{gather*}
\phi(s, t):=H(2 H-1)|t-s|^{2 H-2} ,\\
\|f\|_\phi^2  :=\int_0^{\infty} \int_0^{\infty} f(s) f(t) \phi(s, t) \,ds\,dt, \\
\langle f, g\rangle_\phi  :=\int_0^{\infty} \int_0^{\infty} f(s) g(t) \phi(s, t) \,ds\,dt,
\end{gather*}
where
$f, g:\mathbb{R} \to \mathbb{R}$ are Borel measurable functions, and define
\[
 \mathbb{L}_\phi^2\left(\mathbb{R}\right)=\left\{f:\mathbb{R} \to \mathbb{R},
 f  \text{ is Borel measurable, } \|f\|_\phi^2<+\infty\right\}.
\]
By considering the step functions
\[
 f_m(t)=\sum_i a_i^{(m)} \chi_{\left[t_i, t_{i+1}\right)}\left(t\right)
\]
approximating a deterministic $f \in \mathbb{L}_\phi^2\left(\mathbb{R}\right)$, define
\[
 \int_{0}^{\infty} f(t) d B_t^H
 =\lim _{m \to \infty} \int_{0}^{\infty}f_m(t) d B_t^H \\
 =\lim _{m \to \infty}\Big(\sum_i a_i^{(m)}\left(B_{t_{i+1}}^H-B_{t_i}^H\right)\Big).
\]

Define the Wick exponentials
\[
 \varepsilon(f):=\exp \Big(\int_0^{\infty} f(t) d B^H_t-\frac{1}{2}\|f\|_\phi^2\Big),
\quad f \in \mathbb{L}_\phi^2\left(\mathbb{R}\right).
\]
The Wick product of two exponentials $\varepsilon(f)$ and $\varepsilon(g)$ 
is defined as
\[
 \varepsilon(f) \diamond \varepsilon(g):=\varepsilon(f+g).
\]

According to \cite[Theorem 3.1]{duncan2000stochastic}, the linear span of 
$\left\{\varepsilon(f), \ f \in \mathbb{L}_\phi^2\left(\mathbb{R}\right)\right\}$ 
is a dense set of $L^2\left( \Omega,\mathcal{F} ,P \right)$. 
In other words, the random variables in $L^2\left( \Omega,\mathcal{F} ,P \right)$ 
can be approximated by Wick exponentials' linear combination and the definition 
of Wick product can be extended to random variables $X, Y \in L^2$ through the 
fractional Wiener-It\^o chaos expansion (see, \cite{hu2003fractional}). 
Let $X_s \in L^2\left( \Omega,\mathcal{F}_s ,P \right)$, $s \in \left[0,T\right]$. 
If the Wick product and the limit exists, the fractional Wick-It\^o integral with 
respect to fractional Brownian motion $\left(B_s^H\right)_{[0, T]}$ can be defined 
by Wick-Riemann sums such as
\[
 \int_0^T X_s d^{\diamond} B_s^H:=\lim _{n \to \infty} 
 \sum_{t_i \in \pi_n} X_{t_{i-1}} \diamond\big(B_{t_i}^H-B_{t_{i-1}}^H\big),
\]
where $\pi_n=\left\{0=t_0<t_1<\dots<t_n=T\right\}$ with
$\max _{t_i \in \pi_n}\left|t_i-t_{i-1}\right| \to 0$ for $n \to \infty$.

Based on the work of \cite{sottinen2001fractional} who approximated fractional 
Brownian motion by disturbed binary random walks and the discrete Wick product 
introduced by \cite{holden1992discrete}, \cite{bender2004arbitrage} proposed the 
discrete Wick product on fractional Brownian motion to avoid the technical difficulties 
in computing the continuous Wick product.

We define the discrete Wick exponential 
$\exp ^{\diamond_n}\left(I^n\left(f^n\right)\right)$, that is,
\begin{gather*}
 I^n\left(f^n\right):=\frac{1}{\sqrt{n}} \sum_{i=1}^n f_i^n \xi_i^n, \\
 \exp ^{\diamond_n}\left(I^n\left(f^n\right)\right)
 :=\prod_{i=0}^n\Big(1+\frac{1}{\sqrt{n}} f_i^n \xi_i^n\Big) ,
\end{gather*}
where
$f^n=\left(f_i^n, \dots, f_n^n\right) \in \mathbb{R}^n$,
$\left(\xi_1^n, \dots, \xi_n^n\right)$ is an $n$-tuple of independent symmetric
Bernoulli random variables with 
$P_n\left(\xi_i^n=1\right)=P_n\left(\xi_i^n=-1\right)=\frac{1}{2}$, 
living on a probability space $\left(\Omega_n, \mathcal{F}_n, P_n\right)$.
Analogously, a discrete Wick product of the discrete Wick exponentials 
$I^n\left(f^n\right), I^n\left(g^n\right)$ is
\[
 \exp ^{\diamond_n}\left(I^n\left(f^n\right)\right) \diamond_n 
 \exp ^{\diamond_n}\left(I^n\left(g^n\right)\right)
 =\exp ^{\diamond_n}\left(I^n\left(f^n+g^n\right)\right) .
\]
Obviously, $L^2\left(\Omega_n, \mathcal{F}_n, P_n\right)$ is a $2^n$-dimensional 
vector space, and a canonical orthonormal basis of 
$L^2\left(\Omega_n, \mathcal{F}_n, P_n\right)$ consists of
\[
 \Xi_A^n:=\prod_{i \in A} \xi_i^n, \quad A \subset\{1, \dots, n\}.
\]
Every $X,Y \in L^2\left(\Omega_n, \mathcal{F}_n, P_n\right)$ has a unique expansion 
in terms of this basis, which is called the Walsh decomposition of $X$,
\[
 X=\sum_{A \subset\{1, \dots, n\}} X_A^n \Xi_A^n, \quad
 Y=\sum_{B \subset\{1, \dots, n\}} Y_B^n \Xi_B^n,
\]
where $X_A^n,Y_B^n \in \mathbb{R}$.

Motivated by analogies of the chaos decomposition, the definition of discrete 
Wick product can be extended to random variables 
$X, Y \in L^2\left(\Omega_n, \mathcal{F}_n, P_n\right)$ (see \cite{bender2012connection}),
\[
 X \diamond_n Y=\sum_{C \subset\{1, \dots, n\}}
 \left(\sum\left\{X_A^n Y_B^n: A \cap B=\emptyset, A \cup B=C\right\}\right) \Xi_C^n.
\]

\subsection{Fractional Black-Scholes model and its discrete method}
In fractional Black-Scholes market, the dynamics of the interest rate $r_t$ satisfies
\begin{equation}
\label{fbs}
 d r_t=\mu r_t d t+\sigma r_t d^{\diamond} B_t^H, \ r_0=s_0,
\end{equation}
where $\mu \in \mathbb{R}$ is the drift rate of the risk interest price process, $\sigma>0$ is the volatility, $s_0$ is the initial price. According to fractional It\^o formula in \cite[Theorem 4.3]{duncan2000stochastic}, the solution of the stochastic differential equation \eqref{fbs} is
\[
 r_t=r_0 \exp \left(\mu t-\frac{1}{2} \sigma^2 t^{2 H}+\sigma B_t^H\right).
\]

We define a discrete version of the fractional Black-Scholes model by
\begin{equation}
\label{dfbs}
 r_i^n=\left(1+\frac{\mu}{n}\right) r_{i-1}^n+\sigma r_{i-1}^n \diamond_n\left(B_{\frac{i}{n}}^{H, n}-B_{\frac{i-1}{n}}^{H, n}\right), \ r_0^n=s_0,
\end{equation}
where
\begin{gather*}
 B_t^{H, n}:=\sum_{i=1}^{\lfloor n t\rfloor} b_{t, i}^n \frac{1}{\sqrt{n}} \xi_i^n,\\
 b_{t, i}^n:=n \int_{\frac{i-1}{n}}^{\frac{i}{n}} z_H
 \Big(\frac{\lfloor n t\rfloor}{n}, s\Big) d s, \\
 z_H(t,s)
  = \begin{cases}C_H s^{H-\frac{1}{2}}(t-s)^{H-\frac{1}{2}},
   & \text{if } s<t, \\ 
   0, & \text{otherwise, }
   \end{cases}
\end{gather*}
$ C_H$ is a normalization constant depending on $H$.

The following theorem shows that $r_{\lfloor n t\rfloor}^n$ obtained by the 
discrete version \eqref{dfbs} weakly converges to $r_t$.

\begin{lemma}[{\cite[Theorem 1.4]{bender2012connection}}]
Suppose $\mu, s_0 \in \mathbb{R}, \sigma>0$. 
Then $\widetilde{r}_t^n:=r_{\lfloor n t\rfloor}^n$, the piecewise constant 
interpolation of the discrete Wick difference equation \eqref{dfbs}, converges weakly
to the interest rate $r$ in the fractional Black-Scholes model, i.e. the solution of 
the SDE \eqref{fbs}, in the Skorokhod space $D([0,1], \mathbb{R})$.
\end{lemma}

Equation \eqref{dfbs} shows that recursively functions 
$r_{i}^n\left(x_1,\dots, x_i\right)$ can be constructed such that
$r_{i}^n=r_{i}^n\left(\xi_1^n,\dots, \xi_i^n\right)$.
Applying \cite[Lemma 1.3]{bender2012connection}, the discrete Wick difference 
equation \eqref{dfbs} can be reformulated as
\begin{equation}
\label{wffbs}
 r_i^n=r_{i-1}^n\left(1+\frac{\mu}{n}+\sigma\left(B_{\frac{i}{n}}^{H, n}
 -B_{\frac{i-1}{n}}^{H, n}\right)\right)-\sigma \frac{1}{n} \sum_{j=1}^{i-1}
 \left(b_{\frac{i}{n}, j}^n-b_{\frac{i-1}{n}, j}^n\right) D_j^n r_{i-1}^n ,
\end{equation}
where
\begin{equation}
 D_i^n r_{i-1}^n=\frac{r_{i-1}^n\left(\xi_1^n, \dots,
 \xi_{i-1}^n, 1, \dots, \xi_n^n\right)-r_{i-1}^n\left(\xi_1^n,
 \dots, \xi_{i-1}^n,-1, \dots, \xi_n^n\right)}{2 / \sqrt{n}}
\end{equation}
is discrete Malliavin derivative with respect to the increments 
$\frac{1}{\sqrt{n}} \xi_i, i=1, \dots, n$,
\[
r_i^n=r_{i-1}^n\left(1+\frac{\mu}{n}+\sigma
\left(B_{\frac{i}{n}}^{H, n}-B_{\frac{i-1}{n}}^{H, n}\right)\right).
\]

\section{Main results}
\label{sec:main}

\subsection{Lower bound and confident intervals}
Below are the detailed steps to estimate the lower bound of an interest rate under 
fractional Black-Scholes model.
\smallskip

\noindent\textbf{Step 1.} Simulate $n$-tuple of independent symmetric Bernoulli
 random variables
\begin{gather*}
 \xi_0(k_1)=\left(\xi_0^1(k_1), \dots,
  \xi_0^n(k_1)\right), \ k_1=1,\dots,m,\\
 \xi_1\left(k_1,k_2\right)=\left(\xi_1^1\left(k_1,k_2\right), \dots,
 \xi_1^n\left(k_1,k_2\right)\right), \ k_2=1,\dots,m,\\
 \vdots\\
\begin{aligned}
\xi_{N-1}\left(k_1,\dots,k_{N-1}\right)
&=\left(\xi_{N-1}^1\left(k_1,\dots,k_{N-1}\right), \dots, \xi_{N-1}^n\left(k_1,
\dots,k_{N-1}\right)\right), \\ &k_{N-1}=1,\dots,m,
\end{aligned}
\end{gather*}
with
\[ 
P\left(\xi_i^j\left(\cdot\right)=1\right)
=P\left(\xi_i^j\left(\cdot\right)=-1\right)
=\frac{1}{2},\ i=0,\dots,N-1,\quad  j=1,\dots,n.
\]
\smallskip

\noindent\textbf{Step 2.} For each epoch $i$ and $j$, compute $b_{i,p}^{j,q}$ by
\[
 b_{i,p}^{j,q}=n \int_{t_p+\frac{q-1}{n}}^{t_p+\frac{q}{n}} z_H
  \Big(t_{i}+\frac{jT}{nN}, s\Big) d s,\quad  p=0,\dots,i,\ q=1,\dots,n.
\]
\smallskip

\noindent\textbf{Step 3.} For $i=0$, compute $B_{0}^{j}(k_1)$ and
$r_{0}^{j}(k_1)$ by
\begin{gather*}
 B_{0}^{j}(k_1)
 =\sum_{q=1}^{j} b_{0,0}^{j,q} \frac{1}{\sqrt{n}} \xi_0^q(k_1),\\
 r_{0}^{j}(k_1)
 = r_{0}^{j-1}(k_1)\Big(1+\frac{\mu T}{nN}+\sigma\left(B_{0}^{j}(k_1)-B_{0}^{j-1}(k_1)
 \right)\Big)
 - \frac{\sigma T}{nN} \sum_{q=1}^{j-1}\left(b_{0,0}^{j,q}-b_{0,0}^{j,q-1}\right)
  D_q^n r_{0}^{j-1}(k_1) ,
\\
 B_{0}^{0}(k_1)=0,\quad r_{0}^{0}(k_1)=s_0.
\end{gather*}
\smallskip

\noindent\textbf{Step 4.} For $1\le i \le N-1$, compute 
$B_{i}^{j}\left(k_0,\dots,k_i\right)$ and $r_{i}^{j}\left(k_0,\dots,k_i\right)$ by
\begin{equation}
\label{pathgeneration}
\begin{gathered}
 B_{i}^{j}\left(k_0,\dots,k_i\right)
 =\sum_{p=0}^{i-1}\sum_{q=1}^{n}b_{i,p}^{j,q} \frac{\xi_p^q\left(k_0,\dots,k_p\right)}
 {\sqrt{\left(i+1\right)n}} +\sum_{q=1}^{j}b_{i,i}^{j,q} 
 \frac{\xi_i^q\left(k_0,\dots,k_i\right)}{\sqrt{\left(i+1\right)n}} ,\\
\begin{aligned}
r_{i}^{j}\left(k_0,\dots,k_i\right)
&= r_{i}^{j-1}\left(k_0,\dots,k_i\right)\left(1+\frac{\mu T}{nN}+\sigma\left(B_{i}^{j}
 \left(k_0,\dots,k_i\right)-B_{i}^{j-1}\left(k_0,\dots,k_i\right)\right)\right)\\
&\quad - \frac{\sigma T}{nN} \sum_{q=1}^{j-1}\left(b_{i,i}^{j,q}-b_{i,i}^{j,q-1}\right)
  D_{i-1}^q r_{i}^{j-1}\left(k_0,\dots,k_i\right) ,
\end{aligned}\\
 B_{i}^{0}\left(k_0,\dots,k_i\right)=B_{i-1}^{n}\left(k_0,\dots,k_{i-1}\right),\\
 r_{i}^{0}\left(k_0,\dots,k_i\right)=r_{i-1}^{n}\left(k_0,\dots,k_{i-1}\right).
\end{gathered}
\end{equation}
For $i=N$, let
\[
 r_{N}^{0}\left(k_0,\dots,k_{i-1}\right)=r_{N-1}^{n}\left(k_0,\dots,k_{i-1}\right).
\]
\smallskip

\noindent\textbf{Step 5.} For each path and epoch, compute the call value of the 
zero-coupon bond,
\[
 h_i\left(r_{i-1}^{n}\left(k_0,\dots,k_{i-1}\right)\right)=
 p_{t_{i-1}}\left(r_{i-1}^n\left(k_0,\dots,k_{i-1}\right)\right)
\]
\smallskip

\noindent\textbf{Step 6.} For each path, compute the continuation value $C_i^*$ of 
the bond, defined as the present value of the expected one-period-ahead bond value,
\[
 C_i^*\left(r_{i-1}^n\left(k_0, \dots, k_{i-1}\right)\right)=\frac{1}{m} \sum_{l=1}^m
 e^{- r_iT/N}
V_{i+1}^*\left(r_i^n\left(k_0, \dots, k_{i-1}, l\right)\right) ,
\]
where $\Delta t_i=t_{i+1}-t_i$ and $V_i^*$ is defined in Step 8 below,
\begin{gather*}
 V_N^*\left(r_{N-1}^{n}\left(k_0,\dots,k_{N-1}\right)\right)
 = h_N\left(r_{N-1}^{n}\left(k_0,\dots,k_{N-1}\right)\right),\\
 C_N^*\left(r_{N-1}^{n}\left(k_0,\dots,k_{N-1}\right)\right)= 0.
\end{gather*}
\smallskip

\noindent\textbf{Step 7.} For each path, define the tentative redeem-or-issue indicator 
variable $x_i\left(k_0,\dots,k_{i-1}\right)$,
\[
 x_i\left(k_0,\dots,k_{i-1}\right)
 =\begin{cases}
1, & \text{if } h_i\left(r_{i-1}^{n}\left(k_0,\dots,k_{i-1}\right)\right)
 \geq C_i^*\left(r_{i-1}^{n}\left(k_0,\dots,k_{i-1}\right)\right) ,\\
0, & \text{if } h_i\left(r_{i-1}^{n}\left(k_0,\dots,k_{i-1}\right)
 \right)<C_i^*\left(r_{i-1}^{n}\left(k_0,\dots,k_{i-1}\right)\right).
\end{cases}
\]
\smallskip

\noindent\textbf{Step 8.} For each path $k$, define the current value 
$V_i^*\left(k_0,\dots,k_{i-1}\right)$ of the bond
\[
 V_i^*\left(k_0,\dots,k_{i-1}\right)
 = \begin{cases}
 h_i\left(k_0,\dots,k_{i-1}\right) & \text{if } x_i\left(k_0,\dots,k_{i-1}\right)=1, \\
 C_i^*\left(k_0,\dots,k_{i-1}\right) & \text{if } x_i\left(k_0,\dots,k_{i-1}\right)=0 .
 \end{cases}
\]
\smallskip

\noindent\textbf{Step 9.} If $i>0$, then set $t_i=t_{i-1}$ and return to Step 6, 
otherwise stop the iterations and compute the redeem-or-issue indicator variable 
$y_i\left(k_0,\dots,k_{i-1}\right)$ and optimal stopping time
$\hat{\tau}\left(k_0,\dots,k_{N-1}\right)$
\begin{gather*}
 y_i\left(k_0,\dots,k_{i-1}\right)
 = \begin{cases}1 & \text{if } x_i\left(k_0,\dots,k_{i-1}\right)=1 
 \text{ and } x_j\left(k_0,\dots,k_{j-1}\right)=0 \text{ for all } j<i, \\
 0 & \text{otherwise. }
 \end{cases}
\\
 \hat{\tau}\left(k_0,\dots,k_{N-1}\right)
 = \begin{cases}i & \text{if } y_i\left(k_0,\dots,k_{i-1}\right)=1 , \\
 N & \text{otherwise.}
 \end{cases}
\end{gather*}
\smallskip

\noindent\textbf{Step 10.} The price of the zero-coupon callable bond is estimated by
\begin{equation}
\label{hatv}
 \hat{V}=\frac{1}{m^N} \sum_{k_0=0}^m \sum_{k_1=0}^m \dots \sum_{k_N-1=0}^m
 \sum_{i=0}^N e^{-r_{i-1} t_i} y_i\left(k_0,\dots,k_{i-1}\right)
  h_i\left(r_{i-1}^{n}\left(k_0,\dots,k_{i-1}\right)\right).
\end{equation}

For  establishing profs of convergence, we use the following assumptions.
\begin{itemize}  %\label{Ass1}
\item[(A1)] For $i \in\{0, \dots, N-1\}$, we have
\[
 E[h_i^2]<\infty, \text{ and } E[\left(r_{i-1}^n\right)^2]<\infty .
\]
\item[(A2)] %\label{Ass2}
 $P[C_i=h_i]=0$ for $i=1, \dots, N-1$.
\end{itemize}
Let $F_i\left(x\mid \mathcal{F}_i \right)$ be the conditional probability of 
$V_{i+1}\left(r_{i+1}^0\right)$ given $\xi_0,\dots,\xi_{i-1}$, that is,
\[
 F_i\left(x\mid \mathcal{F}_i \right)
 =P[V_{i+1}\left(r_{i}^n\right) \leq x \mid \xi_0,\dots,\xi_{i-1}], \quad
 i=0,1, \dots, N-1 .
\]
Let the estimator for $F_i\left(x\mid \mathcal{F}_i \right)$ be defined by
\[
 \hat{F}_i^m\left(x\mid \mathcal{F}_i \right)
 =\frac{N\left(x, \xi_0,\dots,\xi_{i-1}\right)}{m},
\]
where $N\left(x, \xi_0,\dots,\xi_{i-1}\right)$ is the number of path index such that
\[
 r_{i+1}^{0}\left(k_0,\dots,k_{i-1},l\right) \leq x.
\]

\begin{lemma} \label{FF}
 For $i=0,1, \dots, N-1$, we have that
\[
 \lim _{m \to \infty} P\Big( \sup _{-\infty<x<\infty}
 \big|\hat{F}_i^m\left(x\mid \mathcal{F}_i \right)-F_i\left(x\mid \mathcal{F}_i 
 \right)\big|=0\Big)=1.
\]
\end{lemma}

\begin{proof}
The proof is based on the Glivenko-Cantelli theorem. For any positive integer $c$, 
take $x_{c, k}$ to be the smallest $x$ that satisfies
\[
 F_i\left(x-0\mid \mathcal{F}_i \right)=F_i\left(x\mid \mathcal{F}_i \right) 
 \leq \frac{k}{c} \leq F_i\left(x+0\mid \mathcal{F}_i \right), \quad k=1,2, \dots, c.
\]
By the Law of Large Numbers and \eqref{wffbs},
\begin{gather*}
 P\Big(\lim _{m \to \infty} \hat{F}_i^m(x_{c, k}\mid \mathcal{F}_i )
 =F_i(x_{c, k}\mid \mathcal{F}_i)\Big)=1,
\\
 P\Big(\lim _{m \to \infty}\hat{F}_i^m\left(x_{c, k}+0\mid \mathcal{F}_i \right)
 =F_i\left(x_{c, k}+0\mid \mathcal{F}_i\right)\Big)=1.
\end{gather*}
Let
\begin{gather*}
 \begin{aligned}
A_k^c & =\left\{\lim _{m \to \infty} \hat{F}_i^m\left(x_{c, k}\mid \mathcal{F}_i \right)=F_i\left(x_{c, k}\mid \mathcal{F}_i\right)\right\}\\
&=\left\{\lim _{m \to \infty}\left|\hat{F}_i^m\left(x_{c, k}\mid \mathcal{F}_i 
\right)-F_i\left(x_{c, k}\mid \mathcal{F}_i\right)\right|=0\right\},
\end{aligned} \\
\begin{aligned}
B_k^c & =\left\{\lim _{m \to \infty} \hat{F}_i^m\left(x_{c, k}+0\mid \mathcal{F}_i \right)=F_i\left(x_{c, k}+0\mid \mathcal{F}_i\right)\right\} \\
& =\left\{\lim _{m \to \infty}\left|\hat{F}_i^m\left(x_{c, k}+0\mid \mathcal{F}_i 
\right)-F_i\left(x_{c, k}+0\mid \mathcal{F}_i\right)\right|=0\right\} ,
\end{aligned}\\
\begin{aligned}
A^c & =\cap_{k=1}^c\left(A_k^c \cap B_k^c\right)=\left\{\lim_{ m\to \infty } \max_{ 1 \leq k \leq c } \left\{\max \left(\left|\hat{F}_i^m\left(x_{c, k}\mid \mathcal{F}_i \right)-F_i\left(x_{c, k}\mid \mathcal{F}_i\right)\right|,\right.\right.\right. \\
&\quad \left.\left.\left.\left|\hat{F}_i^m\left(x_{c, k}+0\mid \mathcal{F}_i 
\right)-F_i\left(x_{c, k}+0\mid \mathcal{F}_i\right)\right|\right)\right\}=0\right\} ,
\end{aligned}\\
A =\cap_{c=1}^{\infty} A^c.
\end{gather*}
Obviously, $P\left(A_k^c\right)=P\left(B_k^c\right)=1$, and
\[
P\left(\overline{A^c}\right) 
=P\Big(\cup_{k=1}^c\left(\overline{A_k^c} \cup \overline{B_k^c}\right)\Big) 
\leq \sum_{k=1}^c\left(P\left(\overline{A_k^c}\right)+P\left(\overline{B_k^c}\right)
\right)=0.
\]
Thus,
\[
P( \overline{A}) =P\Big(\cup_{c=1}^{\infty} \overline{A^c}\Big)
=P\Big(\lim _{n \to \infty} \cup_{c=1}^n \overline{A^c}\Big) 
\leq \lim _{n \to \infty} \sum_{c=1}^n P\left(\overline{A^c}\right)=0.
\]
For $k=0,1,2, \dots, c$,
\begin{align}
\hat{F}_i^m\left(x\mid \mathcal{F}_i \right)-F_i\left(x\mid \mathcal{F}_i \right) & \leq \hat{F}_i^m\left(x_{c, k+1}\mid \mathcal{F}_i \right)-F_i\left(x_{c, k}+0\mid \mathcal{F}_i\right) \\
& \leq \max_{ 1 \leq k \leq c }\left|\hat{F}_i^m\left(x_{c, k}\mid \mathcal{F}_i \right)-F_i\left(x_{c, k}\mid \mathcal{F}_i\right)\right|+\frac{1}{c}\\
F_i\left(x\mid \mathcal{F}_i \right)-\hat{F}_i^m\left(x\mid \mathcal{F}_i \right) & \leq F_i\left(x_{c, k+1}\mid \mathcal{F}_i\right)-\hat{F}_i^m\left(x_{c, k}+0\mid \mathcal{F}_i \right) \\
& \leq \frac{1}{c}+\max_{ 1 \leq k \leq c }\left|F_i\left(x_{c, k}
+0\mid \mathcal{F}_i\right)-\hat{F}_i^m\left(x_{c, k}+0\mid \mathcal{F}_i \right)\right|.
\end{align}
This completes the proof by letting $c$ and $m$ approach infinity.
\end{proof}

We define
\[
 \widetilde{C}_i^m\left(r_{i-1}^{n}\left(k_0,\dots,k_{i-1}\right)\right)
 =\frac{e^{-rT/N}}{m} \sum_{l=1}^{m} V_{i+1}\left(r_{i}^{n}
 \left(k_0,\dots,k_{i-1},l\right)\right).
\]

\begin{lemma}
\label{CC1}
Assume that the prices of asset are bounded almost surely. Then we have 
 \[
 \lim _{m \to \infty} E\left[\left|\widetilde{C}_i^m\left(r_{i-1}^n\right)-C_i\left(r_{i-1}^n\right)\right|\right] = 0
 \]
for $i=0,1, \dots, N-1$.
\end{lemma}

\begin{proof}
From \eqref{dynamic}, obviously, $\widetilde{C}_i^m\left(r_{i-1}^n\right)$ and 
$C_i\left(r_{i-1}^n\right)$ are all bounded for $m=1,2, \dots, \ i=1, \dots, N$.
We define
\[
 M=\max \left\{\max_i C_i\left(r_{i-1}^n\right), \; 
 \max _{i, m} \widetilde{C}_i^m\left(r_{i-1}^n\right)\right\} .
\]
From the definition of $\widetilde{C}_i^m$,
 \begin{align*}
E\left[\left|\widetilde{C}_i^m\left(r_{i-1}^n\right)-C_i\left(r_{i-1}^n\right)\right| \mid \mathcal{F}_i\right] & =E\Big[\Big|\int_0^M x d \hat{F}_i^m\left(x\mid \mathcal{F}_i \right)-\int_0^M x d F_i\left(x\mid \mathcal{F}_i \right)\Big|\Big] \\
& \leq M\int_0^M E\Big[\Big|\hat{F}_i^m\left(x\mid \mathcal{F}_i \right)
-F_i\left(x\mid \mathcal{F}_i \right)\Big|\Big] d x .
\end{align*}
By Lemma \ref{FF},
\[
 \lim _{m \to \infty} E\left[\left|\hat{F}_i^m\left(x\mid \mathcal{F}_i \right)-F_i\left(x\mid \mathcal{F}_i \right)\right|\right]=0.
\]
Hence, by dominated convergence theorem, we complete the proof.
\end{proof}

\begin{lemma}
\label{CC2}
Under assumptions of Lemma \ref{CC1}, for $i=0,1, \dots, N-1$,
 \begin{equation}
 \label{l41}
 \lim _{m \to \infty} E\left[\left|C_i^*\left(r_{i-1}^{n}\right)-\widetilde{C}_i^m\left(r_{i-1}^n\right)\right|\right]=0,
 \end{equation}
and therefore,
\begin{equation}
\label{l42}
 \lim _{m \to \infty} E\left[\left|C_i^*\left(r_{i-1}^{n}\right)-C_i\left(r_{i-1}^{n}\right)\right|\right]=0.
\end{equation}
\end{lemma}

\begin{proof}
The proof uses a recursive method. For $i=N-1$, 
$C_N^*\left(r_{N-1}^{n}\right)=\widetilde{C}_N^m\left(r_{N-1}^n\right)
=C_N\left(r_{N-1}^{n}\right)=0$,
\begin{align*}
 C_{N-1}^*\left(r_{N-2}^{n}\right)
 &= \frac{e^{-rT/N}}{m} \sum_{l=1}^{m} V_{N}^*\left(r_{N-1}^{n}\left(k_0,\dots,l\right)\right) \\
 &= \frac{e^{-rT/N}}{m} \sum_{l=1}^{m} \max \left\{h_N\left(r_{N-1}^{n}\left(k_0,\dots,l\right)\right), \ C_N^*\left(r_{N-1}^{n}\left(k_0,\dots,l\right)\right) \right\} \\
 &= \frac{e^{-rT/N}}{m} \sum_{l=1}^{m} \max \left\{h_N\left(r_{N-1}^{n}\left(k_0,\dots,l\right)\right), \ C_N\left(r_{N-1}^{n}\left(k_0,\dots,l\right)\right) \right\} \\
 &= \widetilde{C}_{N-1}^m\left(r_{N-2}^n\right).
\end{align*}
Hence, \eqref{l41} and \eqref{l42} hold for $i=N-1$. For $i=N-2$,
\begin{gather*}
 C_{N-1}^*\left(r_{N-2}^{n}\left(k_0,\dots,l\right)\right), \quad l=1,\dots,m,\\
 C_{N-1}\left(r_{N-2}^{n}\left(k_0,\dots,l\right)\right), \quad l=1,\dots,m,
\end{gather*}
are independent and identically distributed, respectively. Moreover,
\begin{align*}
& E\left[\left|C_{N-2}^*\left(r_{N-3}^{n}\right)-\widetilde{C}_{N-2}^m\left(r_{N-3}^n\right)\right| \mid \mathcal{F}_{N-2}\right] \\
&=  E\Big[\Big| \frac{e^{-rT/N}}{m} \sum_{l=1}^{m} \max \left\{h_{N-1}
 \left(r_{N-2}^{n}\left(k_0,\dots,l\right)\right), \;
  C_{N-1}^*\left(r_{N-2}^{n}\left(k_0,\dots,l\right)\right)\right\} \\
&\quad -\frac{e^{-rT/N}}{m} \sum_{l=1}^{m} \max \left\{h_{N-1}
 \left(r_{N-2}^{n}\left(k_0,\dots,l\right)\right), \;
  C_{N-1}\left(r_{N-2}^{n}\left(k_0,\dots,l\right)\right)\right\}\Big| 
  \mid \mathcal{F}_{N-2} \Big]\\
&\leq \frac{e^{-rT/N}}{m} \sum_{l=1}^{m} 
  E\left[\left| \max \left\{h_{N-1}\left(r_{N-2}^{n}\left(k_0,\dots,l\right)\right), \ C_{N-1}^*\left(r_{N-2}^{n}\left(k_0,\dots,l\right)\right)\right\}\right.\right. \\
&\quad \left.\left.-\max \left\{h_{N-1}\left(r_{N-2}^{n}\left(k_0,\dots,l\right)\right), \ C_{N-1}\left(r_{N-2}^{n}\left(k_0,\dots,l\right)\right)\right\}\right| \mid\mathcal{F}_{N-2}\right] \\
&=  e^{-rT/N} E\left[\left|\max \left\{h_{N-1}\left(r_{N-2}^{n}\right), \ C_{N-1}^*\left(r_{N-2}^{n}\right)\right\}\right.\right. \\
&\quad \left.\left.-\max \left\{h_{N-1}\left(r_{N-2}^{n}\right), \;
  C_{N-1}\left(r_{N-2}^{n}\right)\right\}\right| \mid\mathcal{F}_{N-2}\right]\\
&\leq e^{-rT/N} E\left[\left|C_{N-1}^*\left(r_{N-2}^{n}\right)-C_{N-1}
 \left(r_{N-2}^{n}\right)\right| \mid \mathcal{F}_{N-2}\right] , 
\end{align*}
Taking conditional expectation, we have that
\[
 E\left[\left|C_{N-2}^*\left(r_{N-3}^{n}\right)-\widetilde{C}_{N-2}^m
 \left(r_{N-3}^n\right)\right| \right] \leq e^{-rT/N}
  E\left[\left|C_{N-1}^*\left(r_{N-2}^{n}\right)-C_{N-1}\left(r_{N-2}^{n}
  \right)\right| \right]
\]
Hence, \eqref{l41} and \eqref{l42} hold for $i=N-2$. 
The proof is completed by recurrence.
\end{proof}

\begin{lemma}
\label{CC3}
Under Assumption {\rm (A1)}, without imposing boundedness condition, we have that
\[
\lim _{m \to \infty} E\left[\left|C_i^*\left(r_{i-1}^{n}\right)
-C_i\left(r_{i-1}^{n}\right)\right|\right]=0, \quad i=0,1, \dots, N-1. 
\]
\end{lemma}

\begin{proof}
Define the event
\[
 \mathcal{F}_M=\left\{\left|h_i\left(\hat{r}_{t_i}\right)\right| 
 \leq M,|| \hat{r}_{t_i} \|_2 \leq M, i=1, \dots, N\right\}, \quad M>0 .
\]
From Assumption (A1),
\[
 \lim _{M \to \infty} P\left[\mathcal{F}_M^c\right]=0, 
 \quad \max_i E\left[C_i\left(r_{i-1}^n\right)^2\right]<\infty,\quad
 \max _{i, m} E\left[\widetilde{C}_i^m\left(r_{i-1}^n\right)^2\right].
\]
Hence,
\begin{align*}
&E\left[\left|C_i^*\left(r_{i-1}^{n}\right)-C_i\left(r_{i-1}^{n}\right)\right|\right] \\
&=  P\left[\mathcal{F}_M\right] E\left[\left|C_i^*\left(r_{i-1}^{n}\right)-C_i\left(r_{i-1}^{n}\right)\right| \mathcal{F}_M\right]+E\left[\left|C_i^*\left(r_{i-1}^{n}\right)-C_i\left(r_{i-1}^{n}\right)\right| I\left(\mathcal{F}_M^c\right)\right] \\
&\leq  P\left[\mathcal{F}_M\right] E\left[\left|C_i^*\left(r_{i-1}^{n}\right)-C_i\left(r_{i-1}^{n}\right)\right| \mathcal{F}_M\right]+\left(E\left[\left|C_i^*\left(r_{i-1}^{n}\right)-C_i\left(r_{i-1}^{n}\right)\right|^2\right]\right)^{1/2} P\left[\mathcal{F}_M^c\right]^{1/2}.
\end{align*}
Let $M \to \infty$, and the Lemma \ref{CC3} is proved by Lemma \ref{CC2}.
\end{proof}

\begin{lemma}
\label{tau}
Under Assumptions {\rm (A1)} and {\rm (A2)}, we have that
 \[
 \lim _{m \to \infty} P\left[\hat{\tau}_m \neq \tau\right]=0 .
 \]
\end{lemma}

\begin{proof}
The proof is based on the work of Broadie and Glasserman \cite{broadie2004stochastic}.
\begin{align*}
P\left[\hat{\tau}_m \neq \tau\right]
&=  P\left[\hat{\tau}_m<\tau\right]+P\left[\hat{\tau}_m>\tau\right] \\
&=  P\left\{\exists i, \text{ s.t. } C_i^*\left(r_{i-1}^{n}\right) \leq h_i\left(r_{i-1}^{n}\right)< C_i\left(r_{i-1}^{n}\right)\right\} \\
&\quad +P\left\{\exists i, \text{ s.t. } C_i^*\left(r_{i-1}^{n}\right) 
 > h_i\left(r_{i-1}^{n}\right) \geq C_i\left(r_{i-1}^{n}\right)\right\} \\
&\leq  \sum_{i=0}^{N-1} P\left\{C_i^*\left(r_{i-1}^{n}\right) 
 \leq h_i\left(r_{i-1}^{n}\right)< C_i\left(r_{i-1}^{n}\right)\right\} \\
&\quad +\sum_{i=0}^{N-1} P\left\{C_i^*\left(r_{i-1}^{n}\right) 
 > h_i\left(r_{i-1}^{n}\right) \geq C_i\left(r_{i-1}^{n}\right)\right\} .
\end{align*}
From (A1), for any $\epsilon>0$, there exists $\gamma>0$ such that
\[
 P\left[\left|h_i-C_i\right| \geq \gamma\right] \leq \epsilon.
\]
Thus,
\[
 P\left[\hat{\tau}_m \neq \tau\right] 
 \leq N \epsilon+\sum_{i=0}^{N-1} P\left[\left|C_i^*\left(r_{i-1}^{n}\right)
 -C_i\left(r_{i-1}^{n}\right)\right| \geq \gamma\right].
\]
By Lemma \ref{CC3}, it can be proved that
$ \lim _{m \to \infty} P\left[\hat{\tau}_n \neq \tau\right]=0$,
since $\epsilon$ is an arbitrary positive number.
\end{proof}

\begin{theorem}
 Under Assumptions {\rm (A1)} and {\rm (A2)}, we have that
\[
 \lim _{m \to \infty} E\left[\left|\hat{V}-V_0\right|\right]=0 .
\]
\end{theorem}

\begin{proof}
From (A1), by the Law of Large Numbers and dominated convergence theorem,
\[
 \lim _{m \to \infty} E\Big[\Big| \frac{1}{m^N} 
 \sum_{k_0=0}^m \sum_{k_1=0}^m \dots \sum_{k_N-1=0}^m 
 e^{-r \tau} h_{\tau}\left(r_{\tau-1}^n\right) - V_0\left(r_0 \right) \Big|\Big]=0 .
\]
Therefore, all we need to prove is that
\[
 \lim _{m \to \infty} E\Big[\Big| \frac{1}{m^N} \sum_{k_0=0}^m \sum_{k_1=0}^m 
 \dots \sum_{k_N-1=0}^m e^{-r \tau} h_{\tau}\left(r_{\tau-1}^n\right) 
 - \hat{V}\left(r_0 \right) \Big|\Big]=0 .
\]
According to \eqref{hatv}, $\hat{V}$ can be rewritten as
\[
 \hat{V}=\frac{1}{m^N} \sum_{k_0=0}^m \sum_{k_1=0}^m \dots 
 \sum_{k_N-1=0}^m e^{-r \hat{\tau}} h_{\hat{\tau}}\left(r_{\hat{\tau}-1}^n\right).
\]
Hence,
\begin{align*}
& E\Big[\Big| \frac{1}{m^N} \sum_{k_0=0}^m \sum_{k_1=0}^m \dots 
 \sum_{k_N-1=0}^m e^{-r \tau} h_{\tau}\left(r_{\tau-1}^n\right) 
 - \hat{V}\left(r_0 \right)\Big|\Big] \\
&= E\Big[\Big| \frac{1}{m^N} \sum_{k_0=0}^m \sum_{k_1=0}^m \dots 
 \sum_{k_N-1=0}^m \left(e^{-r \tau} h_{\tau}\left(r_{\tau-1}^n\right)
  -e^{-r \hat{\tau}} h_{\hat{\tau}}\left(r_{\hat{\tau}-1}^n\right)\right) 
  I\left(\hat{\tau} \neq \tau\right)\Big|\Big] \\
&\leq  E\left[\left|\left(e^{-r \tau} h_{\tau}\left(r_{\tau-1}^n\right)
 -e^{-r \hat{\tau}} h_{\hat{\tau}}\left(r_{\hat{\tau}-1}^n\right)\right) 
  I\left(\hat{\tau} \neq \tau\right)\right|\right] \\
&\leq \left(E\left[\left(e^{-r \tau} h_{\tau}\left(r_{\tau-1}^n\right)
 -e^{-r \hat{\tau}} h_{\hat{\tau}}\left(r_{\hat{\tau}-1}^n\right)
  \right)^2\right]\right)^{1 / 2}\left(P\left(\hat{\tau} 
  \neq \tau\right)\right)^{1 / 2} .
\end{align*}
Letting $m \to \infty$, the theorem is proved by Lemma \ref{tau}.
\end{proof}

\begin{remark} \rm
Based on the dual formulation proposed by Rogers \cite{rogers2002monte}, we can 
solve the upper bound estimate for zero-coupon bonds.
Since calculating the upper bound is more complex, we will address it in future 
research by providing the point estimate of the bound along with its 95\% confidence 
interval. For now, we only present the expression for the point estimate.
 The point estimate of $V_0$ is
$ \frac{\hat{L}+\hat{U}}{2}$.
\end{remark}

\begin{thebibliography}{99}
\bibitem[1]{aaron2021zero} N. Aaron Pancost;
Zero-coupon yields and the cross-section of bond prices,
\emph{The Review of Asset Pricing Studies}, \textbf{11}(2) (2021), 209-268.
\bibitem[3]{becker2019deep} Sebastian Becker, Patrick Cheridito, Arnulf Jentzen;
Deep optimal stopping,
\emph{Journal of Machine Learning Research}, \textbf{20}(74) (2019), 1-25.
\bibitem[4]{becker2021solving} Sebastian Becker, Patrick Cheridito, Arnulf Jentzen, Timo Welti;
Solving high-dimensional optimal stopping problems using deep learning,
\emph{European Journal of Applied Mathematics}, \textbf{32}(3) (2021), 470-514.
\bibitem[5]{bender2004arbitrage} Christian Bender, Robert J. Elliott;
Arbitrage in a discrete version of the Wick-fractional Black-Scholes market,
\emph{Mathematics of Operations Research}, \textbf{29}(4) (2004), 935-945.
\bibitem[6]{bender2012connection} Christian Bender, Peter Parczewski;
On the connection between discrete and continuous Wick calculus with an application to the fractional Black-Scholes model,
In \emph{Stochastic Processes, Finance and Control: A Festschrift in Honor of Robert J. Elliott}, (2012), 3-40. World Scientific.
\bibitem[7]{broadie2004stochastic} Mark Broadie, Paul Glasserman;
A stochastic mesh method for pricing high-dimensional american options,
\emph{Journal of Computational Finance}, \textbf{7} (2004), 35-72.
\bibitem[8]{diaz2022yield} Antonio D\'{i}az, Francisco Jare\~{n}o, Eliseo Navarro;
Yield curve data choice and potential moral hazard: An empirical exercise on pricing callable bonds,
\emph{International Journal of Finance $\&$ Economics}, \textbf{27}(2) (2022), 2124-2145.
\bibitem[9]{dow2025interest} Scott S. Dow, Stefanos C. Orfanos;
Interest rate sensitivity of callable bonds and higher-order approximations,
\emph{Risks}, \textbf{13}(4) (2025), 69.
\bibitem[10]{duncan2000stochastic} Tyrone E. Duncan, Yaozhong Hu, Bozenna Pasik-Duncan;
Stochastic calculus for fractional Brownian motion I. theory,
\emph{SIAM Journal on Control and Optimization}, \textbf{38}(2) (2000), 582-612.
\bibitem[11]{farto2005numerical} Javier Farto, Carlos V\'{a}zquez;
Numerical techniques for pricing callable bonds with notice,
\emph{Applied Mathematics and Computation}, \textbf{161}(3) (2005), 989-1013.
\bibitem[12]{hobson2024callable} David Hobson, Gechun Liang, Edward Wang;
Callable convertible bonds under liquidity constraints and hybrid priorities,
\emph{SIAM Journal on Financial Mathematics}, \textbf{15}(4) (2024), 1083-1123.
\bibitem[13]{holden1992discrete} Helge Holden, Tom Lindstr{\o}m, Bernt {\O}ksendal, 
 Jan Ub{\o}e;
Discrete Wick calculus and stochastic functional equations,
\emph{Potential Analysis}, \textbf{1}(3) (1992), 291-306.
\bibitem[14]{hu2003fractional} Yaozhong Hu, Bernt {\O}ksendal;
Fractional white noise calculus and applications to finance,
\emph{Infinite Dimensional Analysis, Quantum Probability and Related Topics}, \textbf{6}(01) (2003), 1-32.
\bibitem[15]{kusnadi2023callable} Felivia Kusnadi, Devina Gabriella Tirtasaputra, et al.;
Callable bond's value analysis using binomial interest rate tree considering early redemption and default risks,
\emph{JTAM (Jurnal Teori dan Aplikasi Matematika)}, \textbf{7}(2) (2023), 283-297.
\bibitem[16]{liang2019dynkin} Gechun Liang, Haodong Sun;
Dynkin games with poisson random intervention times,
\emph{SIAM Journal on Control and Optimization}, \textbf{57}(4) (2019), 2962-2991.
\bibitem[17]{lin2022pricing} Sha Lin, Songping Zhu;
Pricing callable-puttable convertible bonds with an integral equation approach,
\emph{Journal of Futures Markets}, \textbf{42}(10) (2022), 1856-1911.
\bibitem[18]{narayanan1989call} M. P. Narayanan, Suk-Pil Lim;
On the call provision in corporate zero-coupon bonds,
\emph{Journal of Financial and Quantitative Analysis}, \textbf{24}(1) (1989), 91-103.
\bibitem[19]{rogers2002monte} L. C. G. Rogers;
Monte carlo valuation of American options,
\emph{Mathematical Finance}, \textbf{12}(3) (2002), 271-286.
\bibitem[20]{skalicky2022valuation} Roman Skalick`{y}, Marek Zinecker, Adam P. Balcerzak, 
Michal Bernard Pietrzak,  El\.{z}bieta Rogalska;
Valuation of embedded options in non-marketable callable bonds: A new numerical approach,
\emph{Technological and Economic Development of Economy}, \textbf{28}(4) (2022), 1115-1136.
\bibitem[21]{sottinen2001fractional} Tommi Sottinen;
Fractional Brownian motion, random walks and binary market models,
\emph{Finance and Stochastics}, \textbf{5}(3) (2001), 343-355.
\bibitem[22]{tan2022deeppricing} Xiaoyu Tan, Zili Zhang, Xuejun Zhao, Shuyi Wang;
Deeppricing: pricing convertible bonds based on financial time-series generative adversarial networks,
\emph{Financial Innovation}, \textbf{8}(1) (2022), 64.
\end{thebibliography}
\end{document}
